\documentclass[11pt,a4paper]{article}
\usepackage[margin=20mm,headheight=16pt]{geometry}
\usepackage{fontspec,kotex,booktabs,longtable,array,xcolor,hyperref,fancyhdr}
\setmainfont{DejaVu Serif}
\setsansfont{DejaVu Sans}
\setmonofont{DejaVu Sans Mono}
\setmainhangulfont{Noto Sans CJK KR}
\setmonohangulfont{Noto Sans CJK KR}
\definecolor{navy}{HTML}{173E59}
\hypersetup{colorlinks=true,linkcolor=navy,urlcolor=navy,
pdftitle={GuardSynth: Related Work},pdfauthor={GuardSynth project}}
\urlstyle{same}
\setlength{\parindent}{0pt}
\setlength{\parskip}{6pt}
\setlength{\emergencystretch}{4em}
\setlength{\tabcolsep}{4pt}
\renewcommand{\arraystretch}{1.2}
\pagestyle{fancy}\fancyhf{}
\fancyhead[L]{\small GuardSynth}
\fancyhead[R]{\small Related Work · v0.1 · 2026-09-30}
\fancyfoot[C]{\thepage}
\renewcommand{\headrulewidth}{0.3pt}
\clubpenalty=10000\widowpenalty=10000
\raggedbottom
\begin{document}
\section*{II. 관련연구}


\phantomsection\section*{A. 추론 기반 자율주행과 CoC}\addcontentsline{toc}{section}{A. 추론 기반 자율주행과 CoC}

시각 정보를 언어적 추론 및 행동으로 연결하는 연구는 자율주행 학습의 중요한 방향이다. Alpamayo-R1은 Chain of Causation(CoC) 추론과 궤적 예측을 결합하고, 지도 미세조정과 강화학습을 통해 추론과 행동의 정합성을 높인다. DriveLM은 인지·예측·계획에 관한 질의응답의 관계를 그래프로 구성하며, DriveVLM은 장면 기술, 장면 분석, 계층적 계획을 연결한다. 이들 연구는 관찰에서 판단에 이르는 과정을 학습 가능한 표현으로 제공한다. \href{\detokenize{https://arxiv.org/abs/2511.00088}}{1}, \href{\detokenize{https://arxiv.org/abs/2312.14150}}{2}, \href{\detokenize{https://arxiv.org/abs/2402.12289}}{3}\par

설명이 행동 선택에 유용하려면 두 결과의 정합성도 중요하다. Drive-R1은 추론과 계획의 불일치 및 과거 정보에 의존하는 지름길 학습을 다루며, 지도 미세조정 이후 궤적과 메타 행동에 관한 보상을 사용한다. 이러한 접근과 상보적으로, 본 연구는 학습 데이터에 제공되는 행동 제약의 표현에 주목한다. 행동의 이유를 설명하는 CoC와, 해당 상황에서 허용되는 행동을 제한하는 제약을 구분하고 두 정보를 결합한다. 이 구분은 기존 추론 모델의 전반적인 성능 부족을 전제하기보다, 설명과 함께 명시적으로 제공할 규범적 정보를 정의하기 위한 것이다. \href{\detokenize{https://ojs.aaai.org/index.php/AAAI/article/view/37602}}{4}\par

\phantomsection\section*{B. 교통 규칙의 명세와 규정 기반 판단}\addcontentsline{toc}{section}{B. 교통 규칙의 명세와 규정 기반 판단}

교통 규칙을 기계가 처리할 수 있는 표현으로 만드는 연구는 행동 제약의 정형화에 직접적인 기반을 제공한다. Esterle 등은 교통 규칙을 선형 시간 논리(LTL)로 표현하고, 주행 데이터에 대한 규칙 준수 평가에 활용하였다. 이는 규칙의 적용 조건과 시간적 관계를 명시적인 의미론으로 다루는 접근이다. 본 연구도 이러한 원리를 따르되, CoC를 보강하는 제약의 대상·적용·유지·해제 및 근거 연결을 EBLC의 요소로 구조화한다. \href{\detokenize{https://arxiv.org/abs/2007.00330}}{5}\par

규정과 장면을 연결하는 또 다른 방향은 검색 증강 추론이다. DriveReg은 상황에 관련된 교통 규정을 검색하고, 이를 언어 모델의 판단에 연결하여 규정 준수와 해석 가능한 의사결정을 지원한다. 이 연구는 장면에 적합한 규정을 찾아 활용하는 과정을 중심에 둔다. 이에 비해 GuardSynth는 선정된 행동 제약을 명시적인 중간 표현으로 작성하고, 제한된 논리 검증을 거쳐 학습용 자연어로 전달하는 과정에 초점을 둔다. 규정 검색과 제약 명세는 서로 대체하는 기능이 아니라, 근거 확보와 제약 표현이라는 상보적인 역할을 갖는다. \href{\detokenize{https://ojs.aaai.org/index.php/AAAI/article/view/41168}}{6}\par

\phantomsection\section*{C. 자연어–정형 명세 변환과 제어 자연어}\addcontentsline{toc}{section}{C. 자연어–정형 명세 변환과 제어 자연어}

자연어 요구사항을 정형 명세로 변환하는 연구는 사람이 이해하는 표현과 기계가 검증하는 표현을 연결한다. NL2TL은 원자 명제를 추상화한 표현과 언어 모델을 이용하여 자연어를 시간 논리로 변환하고, 영역별 명제 인식과 변환 구조를 분리한다. nl2spec은 자연어 구절과 논리식의 부분 구조를 대응시키고, 사용자가 그 대응을 수정하여 모호성을 해소하도록 지원한다. 전자는 변환 모델의 학습과 영역 적응을, 후자는 대화형 명세 작성과 사용자 수정을 강조한다. \href{\detokenize{https://aclanthology.org/2023.emnlp-main.985/}}{7}, \href{\detokenize{https://arxiv.org/abs/2303.04864}}{8}\par

명세를 사람이 읽을 수 있는 언어로 제공하는 방향에서는 제어 자연어(Controlled Natural Language, CNL)가 중요하다. Attempto Controlled English(ACE)는 문법과 해석 규칙을 제한하여 자연어 형태의 문장과 정형 표현을 연결한다. GuardSynth의 CNL은 ACE를 구현한 언어가 아니라, 선택한 EBLC 프로파일의 필드와 조항을 제한된 문장 패턴에 대응시키는 표현이다. 본 연구는 자유로운 자연어 전체의 의미 동등성보다, 정의된 제약 요소와 생성 문장 사이의 추적 가능한 대응을 다룬다. 따라서 논리 검증, 문장 대응 확인, 장면 근거의 타당성 판단은 구별되는 검토 대상이다. \href{\detokenize{https://attempto.ifi.uzh.ch/site/pubs/papers/FLAIRS0601FuchsN.pdf}}{9}\par

\phantomsection\section*{D. 본 연구의 위치}\addcontentsline{toc}{section}{D. 본 연구의 위치}

위 연구들은 주행 추론의 학습, 교통 규칙의 정형화, 언어와 논리의 연결이라는 기반을 제공한다. GuardSynth는 이들을 바탕으로 CoC 행동 제약의 범위·특성·구성 요소를 정의하고, EBLC 명세와 제한된 논리 검증, CNL 생성, CoC 학습 데이터 보강을 연결한다. 핵심은 새로운 주행 모델이나 범용 번역기를 제안하는 것보다, 학습에 제공할 제약을 일관된 표현과 검증 절차로 구성하는 데 있다. 표 1은 각 연구의 중심 목적을 비교한 것으로, 구현 기능의 전수 조사나 성능 순위표가 아니다.\par

표 1. 관련연구의 중심 목적 및 GuardSynth와의 관계.\par

\begingroup\small
\begin{longtable}{@{}>{\raggedright\arraybackslash}p{\dimexpr 0.2200\linewidth-5.280pt\relax}>{\raggedright\arraybackslash}p{\dimexpr 0.2000\linewidth-4.800pt\relax}>{\raggedright\arraybackslash}p{\dimexpr 0.3000\linewidth-7.200pt\relax}>{\raggedright\arraybackslash}p{\dimexpr 0.2800\linewidth-6.720pt\relax}@{}}
\toprule
\textbf{연구} & \textbf{중심 목적} & \textbf{주요 표현·접근} & \textbf{본 연구와의 연결점} \\
\midrule\endfirsthead
\toprule
\textbf{연구} & \textbf{중심 목적} & \textbf{주요 표현·접근} & \textbf{본 연구와의 연결점} \\
\midrule\endhead
\bottomrule\endfoot
Alpamayo-R1; Drive-R1 [1, 4] & 추론과 주행 행동의 정합성 & CoC/CoT, 지도 학습 및 강화학습 & 행동 선택에 유용한 학습 정보 \\
DriveLM; DriveVLM [2, 3] & 장면에서 계획으로 이어지는 추론 & 그래프 질의응답, 계층적 계획 & 상황·행동 설명의 구성 \\
Esterle 등 [5] & 기계 해석 가능한 교통 규칙 & LTL 규칙 명세 & 조건과 시간 관계의 명시적 의미론 \\
DriveReg [6] & 규정에 근거한 주행 판단 & 규정 검색과 언어 모델 추론 & 적용할 규정·근거의 연결 \\
NL2TL; nl2spec [7, 8] & 자연어 요구사항의 정형화 & 시간 논리 변환, 대화형 수정 & 자연어와 명세 요소의 대응 \\
ACE [9] & 명확하게 해석 가능한 언어 & 제한된 문법과 정형 해석 & 제어된 자연어 표현의 원리 \\
GuardSynth & CoC 행동 제약의 구조화와 학습 활용 & EBLC → 논리 검증 → CNL → CoC 보강 & 제약 모델·검증·학습 입력의 연결 \\
\end{longtable}
\endgroup

이 방법의 학습 활용은 동일한 모델과 학습 예산에서 CoC만 제공하는 P0와 제약을 함께 제공하는 P1·P2를 비교하여 평가한다. P1은 우리의 기존 자연어 제약 방식이며, P2는 EBLC에서 생성한 CNL 방식이다. 따라서 P0 대비 P1·P2가 주 비교이고, P1–P2는 제안 계열 내 표현 방식의 비교이다. 이 실험은 제약을 제공하는 효과를 평가하며, 그 차이를 논리 검증만의 인과적 효과로 해석하지 않는다. 본 논문의 범위는 제약의 명세·검증·CNL 생성과 학습 활용이며, EBLC 기반 테스트 생성이나 런타임 감시는 후속 활용에 해당한다.\par

\phantomsection\section*{참고문헌}\addcontentsline{toc}{section}{참고문헌}

\begin{itemize}
\item [1] NVIDIA, Alpamayo-R1, 2025, 개정 2026. \href{\detokenize{https://arxiv.org/abs/2511.00088}}{원 논문}.
\item [2] C. Sima et al., DriveLM, ECCV, 2024. \href{\detokenize{https://arxiv.org/abs/2312.14150}}{원 논문}.
\item [3] X. Tian et al., DriveVLM, 2024. \href{\detokenize{https://arxiv.org/abs/2402.12289}}{원 논문}.
\item [4] Y. Li et al., Drive-R1, AAAI, vol. 40, no. 8, pp. 6708–6716, 2026. \href{\detokenize{https://ojs.aaai.org/index.php/AAAI/article/view/37602}}{출판 논문}.
\item [5] K. Esterle, L. Gressenbuch, and A. Knoll, 교통 규칙의 기계 해석을 위한 정형화, IEEE CAVS, 2020. \href{\detokenize{https://arxiv.org/abs/2007.00330}}{원 논문}.
\item [6] T. Cai et al., DriveReg, AAAI, vol. 40, no. 45, pp. 38287–38295, 2026. \href{\detokenize{https://ojs.aaai.org/index.php/AAAI/article/view/41168}}{출판 논문}.
\item [7] Y. Chen, R. Gandhi, Y. Zhang, and C. Fan, NL2TL, EMNLP, pp. 15880–15903, 2023. \href{\detokenize{https://aclanthology.org/2023.emnlp-main.985/}}{출판 논문}.
\item [8] M. Cosler, C. Hahn, D. Mendoza, F. Schmitt, and C. Trippel, nl2spec, CAV, 2023. \href{\detokenize{https://arxiv.org/abs/2303.04864}}{원 논문}.
\item [9] N. E. Fuchs, K. Kaljurand, and G. Schneider, Attempto Controlled English의 지식 표현·추론·상호운용·사용자 인터페이스 접근, 2006. \href{\detokenize{https://attempto.ifi.uzh.ch/site/pubs/papers/FLAIRS0601FuchsN.pdf}}{원 논문}.
\end{itemize}
\end{document}